\subsection{Axis norms from a mixed degree-six by degree-three seed}
\label{sec:DE-A3-mixed-axis-norms}

\begin{lemma}[Nodal images for an axis kernel]
\label{lem:DE-A3-axis-nodal-norm}
\status{Proved; complete independent proof audits by root and
reduction-descent passed.}
Let $B=E_1\times E_2$ over an algebraically closed field of
characteristic two, and put
\[
             N=\OO_{E_1}(6O)\boxtimes\OO_{E_2}(3O).
\]
Let $\phi:B\to A$ be a finite etale isogeny of degree $n$ whose
kernel is contained in one of the two factors.
A general member $C\in|N|$ is smooth and maps birationally to
an integral curve $Z=\phi(C)$ with exactly $18(n-1)$ ordinary
nodes and no other singularities.  Thus its normalization is $C$.
Every nonempty open condition on smooth members of $|N|$ can
be imposed at the same time.

If $\phi=\phi_1\times1$, then the norm class is $(6,3n)$.
If $\phi=1\times\phi_2$, then the norm class is $(6n,3)$.
These two formulae do not interchange the two factors.
\end{lemma}
\begin{proof}
The complete system has dimension $6\cdot3-1=17$ and is very
ample.  Its smooth-member locus is therefore a nonempty open.
Write $V=H^0(E_1,\OO(6O))$ and
$W=H^0(E_2,\OO(3O))$.
Fix a nonzero kernel translation $j$.
We verify directly that the condition that $C$ and $C-j$ have
a nontransverse intersection has incidence dimension at most
sixteen in $B\times|N|$.

First suppose that $j$ acts on $E_2$.
At a point $(u,v)$ the two value conditions are independent,
because degree three separates the distinct points $v,v+j$.
Their common coefficient fibre is a projective space of
dimension fifteen.
The value and first derivative at $u$ on the degree-six factor
can be prescribed independently.  A section, to first order
in the first-factor coordinate at $u$, can consequently be
written
\[
                        w_0+t w_1,
                    \qquad w_0,w_1\in W
\]
with $w_0,w_1$ independently arbitrary.
The two value conditions require $w_0(v)=w_0(v+j)=0$.
The resulting space of $w_0$ has dimension one, generated by
the line section through these two points in the plane-cubic
embedding.  Let $b_0,b_1$ be its derivatives at the two points.
They cannot both vanish: a nonzero line section of degree three
cannot have zeros of order at least two at two distinct points.
Meanwhile the two values $a_0=w_1(v)$ and $a_1=w_1(v+j)$ are
independently arbitrary.
The intersection determinant is a nonzero scalar multiple of
\[
                              a_0b_1-a_1b_0
\]
times the free coefficient of $w_0$.
It is therefore a genuine nonzero polynomial on every coefficient
fibre.  Its vanishing cuts that fibre in dimension at most fourteen.
Adding the two-dimensional point base gives the required bound.
This argument includes a chord tangent at either of its two points.

If $j$ acts on $E_1$, evaluation of degree-six sections on the
two first jets at $u,u+j$ is surjective: the kernel bundle has
degree $6-4=2>0$ and hence has vanishing $H^1$.
Evaluation of the degree-three factor on its first jet at $v$
is likewise surjective.  Tensoring these evaluations shows that
the two values and all four first partial derivatives can be
prescribed independently.  After imposing the values, the
tangent determinant is again a nonzero polynomial on every
coefficient fibre.  The same dimension bound follows.
Since the incidence is proper, a general section has only
transverse intersections with each of its nontrivial translates.

We also exclude three preimages of a quotient point.
For distinct $0,j,k$ in the kernel, when the kernel acts on
$E_1$, the three values of a degree-six section are independent
at every three distinct points.  Thus the three-value incidence
has dimension $2+(17-3)=16$.
When the kernel acts on $E_2$, degree-three evaluation at the
three distinct points $v,v+j,v+k$ has rank three except when
\[
                              3v+j+k=0.
\]
This is a finite locus in $E_2$, since multiplication by three
is finite etale in characteristic two.  On that locus the
evaluation rank is exactly two, and otherwise it is three.
The exceptional incidence dimension is therefore
$1+(17-2)=16$, and the generic one is again sixteen.
All these proper incidences miss a nonempty open of $|N|$.
There are only finitely many differences and triples, so their
opens can be intersected.

Choose a smooth member in the resulting open.
No nontrivial translation stabilizes it, since its intersection
with every translate is finite.  Hence the restriction of the
quotient is generically injective.  It is an immersion at every
point because $\phi$ is etale.  There are at most two preimages
of every image point, and when there are two their tangent
directions are distinct.  This proves the asserted nodal
description and birationality.
The source has genus $19$ and the image has arithmetic genus
$18n+1$, so the number of nodes is $18(n-1)$.
Equivalently, summing the intersection number
$N^2=36$ over the $n-1$ nonzero kernel elements counts each
unordered pair twice.  This counting is valid also when the
etale kernel has even order; it is an equality of integers.

For completeness, norm of divisors gives the stated classes.
For $\phi_1\times1$, pullback multiplies the first class
coordinate by $n$, whereas
\[
       \phi^*\operatorname{Nm}_{\phi}(N)
           \simeq\bigotimes_{j\in\ker\phi}t_j^*N
\]
has class $(6n,3n)$.
Thus the norm class is $(6,3n)$.
For $1\times\phi_2$ the same calculation gives $(6n,3)$.
The image divisor is the norm divisor because its generic
degree is one.  Finally all conditions used are nonempty opens
in the irreducible projective space $|N|$, which proves the
last assertion about imposing further open conditions.
\end{proof}

\begin{corollary}[The mixed thin strip with ordinary degree twice an odd number]
\label{cor:DE-A3-mixed-thin-vone}
\status{Proved; complete independent proof and seed-certificate audits
by root and reduction-descent passed.}
Let $E_{\rm o}$ be ordinary and $E_{\rm s}$ supersingular over
an algebraically closed field of characteristic two.
For every odd positive integer $n$, a general smooth member of
\[
          \bigl|\OO_{E_{\rm o}}(6nO)
                     \boxtimes\OO_{E_{\rm s}}(3O)\bigr|
\]
has full stable rank $18n-1$ and relative stable rank $18n-2$.
Consequently the previously excluded degree-one product
systems with $a\equiv2\pmod4$, $b=1$, satisfy
$r_1=9a-2$.
\end{corollary}
\begin{proof}
There is an etale isogeny
$\phi_{\rm s}:E'_{\rm s}\to E_{\rm s}$ of degree $n$:
take a cyclic odd-order subgroup on $E_{\rm s}$ and the dual
of its quotient isogeny.  The source is supersingular by
isogeny invariance.  These are the same precise elliptic
isogeny inputs used in
Lemma~\ref{lem:DE-A3-mixed-odd-factor}.
The degree-six by degree-three seed
Lemma~\ref{lem:root-mixed-six-three-cartier}
gives a nonempty open of smooth curves on
$E_{\rm o}\times E'_{\rm s}$ of genus nineteen and full stable
rank seventeen.
Apply Lemma~\ref{lem:DE-A3-axis-nodal-norm} to
$1\times\phi_{\rm s}$ and intersect its nodal-image open
with that sharp open.
The resulting image has $18(n-1)$ nodes, so its full stable
rank is
\[
                         17+18(n-1)=18n-1.
\]
Here the normalization sequence adds one Frobenius-bijective
graph contribution for each node.
The norm class is $(6n,3)$.
Translation of the target puts it in the exact displayed
line-bundle class: the homomorphism associated to an ample
line bundle is an isogeny and is surjective on points over
the algebraic closure.
Openness of iterated Frobenius rank and the nonempty smooth
open in this very ample system transfer this lower bound
to a general smooth member.
The even ordinary-factor Cartier-kernel bound
Lemma~\ref{lem:DE-A3-mixed-even-bound}
gives the matching relative upper bound $18n-2$.
Subtracting the ambient stable rank one from the full stable
rank proves both equalities.
\end{proof}

\begin{remark}[The direction and remaining scope of the norm construction]
\label{rem:DE-A3-mixed-axis-scope}
\status{Proved; complete independent proof audits by root and
reduction-descent passed.}
The isogeny in the corollary acts on the supersingular factor,
and its degree is odd.
An attempted transfer by an arbitrary-degree etale isogeny on
the ordinary factor would instead produce class $(6,3n)$.
That proposed route to all classes $(6n,3)$ is invalid and is
not used here.  In particular this argument makes no assertion
about the still smaller thin-strip cases with $4\mid a$.
The universal seed and both retained programs were independently
read in full before use; their exact determinant certificate is
$c^5$, valid for every ordinary parameter $c\ne0$.
The dimension bounds and the two norm-class formulae above
hold uniformly, without an inference from a finite collection
of kernels or elliptic curves.
\end{remark}
